\makeatletter
\def\input@path{{./}{../}{../../}{../../../}{../../../../}{preamble/}{../preamble/}{../../preamble/}{../../../preamble/}{book2026-2027/preamble/}}
\makeatother
\input{preamble_exams.tex}
\linespread{1.0}
\setlength{\parskip}{3pt}
\newgeometry{margin=0.55in,top=0.45in,bottom=0.55in}
\hypersetup{pdftitle={Grade 11 Term 1 -- C10 Practice: Choosing a Statistical Model}}
% Compact local headings and title; keep the shared preamble unchanged.
\makeatletter
\renewcommand{\section}{\@startsection{section}{1}{0pt}{7pt}{4pt}{\large\bfseries}}
\renewcommand{\@maketitle}{\begin{center}
  {\Large\bfseries \@title\par}\vspace{5pt}
  {\@author\par}
\end{center}}
\makeatother

\title{Grade 11 -- Term 1\\Choosing a Statistical Model: C10 Practice}
\author{\small GLOBAL ECONOMICS -- 2026--2027\\\small Teacher: Nicol\'as L\'opez Cu\'ellar}
\date{}

\begin{document}
\small
\maketitle

\section*{Your Assignment}
Investigate \textbf{one economic, financial, or operational question} using \textbf{one real dataset} and \textbf{one continuous probability model}. Write your report in \textbf{LaTeX using Overleaf} and submit the compiled PDF through \textbf{Open}.

\textbf{Learning objective (C10).} Fit a continuous model to real observations, calculate quantiles and contextual probabilities, compare model probabilities with empirical relative frequencies, and communicate assumptions and limitations. Justify the model using the data and context, and show mathematical procedures even when using software or AI.

\textbf{During the exam:} choose one of the three models on page 2 and describe plausible real-life data in 1--2 sentences (C10: 6 points). No data collection or calculations are required in that brief response. This guide specifies the separate online report for your selected option; its scope is the investigation below.

\section*{1. Choose a question and obtain data}
State a simple question about a \textbf{continuous variable} and obtain at least \textbf{30 usable numerical observations} from one credible source. Choose data reasonably compatible with your proposed model. Use public, non-sensitive observations, not invented or simulated data or a completed practice problem.

Define $X$, its units, context and target population, observation period, source, final sample size $n$, and how observations were selected. Explain removed observations, missing values, and any transformations (including resulting units). Do not remove observations simply to improve the fit.

\textbf{Data access.} Include a clickable hyperlink to the \textbf{original data source}. If you download a file or create a cleaned dataset, upload the \textbf{exact dataset used} to Google Drive or OneDrive and include its clickable link too. Set viewing access so the teacher can open it without requesting permission. Cite the organization/author, dataset title, date or version where available, and access date.

\textbf{Possible questions} (find and assess a real source before committing):
\begin{itemize}
\item How likely is a shipment's measured mass to exceed a stated handling limit? Look for public logistics or quality-control measurements; consider a normal model if roughly symmetric and bell-shaped.
\item How likely is the wait between service arrivals to exceed a stated time? Look for public arrival timestamps and compute consecutive time gaps within a comparable operating period; consider an exponential model.
\item How likely is a bounded service duration to exceed a stated target? Look for public operational records with a roughly increasing or decreasing histogram; consider a linear density on a finite interval.
\end{itemize}

\section*{2. Summarise the data}
Calculate the sample mean, sample variance using $n-1$, and sample standard deviation from the actual observations:
\[
\bar{x}=\frac{1}{n}\sum_{i=1}^{n}x_i,\qquad
s^2=\frac{1}{n-1}\sum_{i=1}^{n}(x_i-\bar{x})^2,\qquad
s=\sqrt{s^2}.
\]
Show formulas, numerical substitutions and intermediate totals such as $\sum x_i$ and $\sum(x_i-\bar{x})^2$. Give units (squared units for variance). You need not print every observation in the report. Include \textbf{one histogram} with a title, labelled axes, units and clear class intervals; explain the bins briefly. If overlaying a density, use a density histogram or rescale the curve to frequencies.

\section*{3. Select and fit one probability density function}
Use \textbf{one} model from page 2. Write its fitted density with numerical parameters and support. Explain briefly why its shape and assumptions are reasonable and identify one limitation. Check its estimated mean and standard deviation against $\bar{x}$ and $s$, and its shape against the histogram. Equality caused by fitting is not independent evidence of a good fit.

\textbf{AI use.} You may ask AI to suggest or estimate parameters from your actual dataset. Name the tool and reproduce the main prompt used; report the estimated parameters, units, and estimation method. Check its calculations against your sample statistics and the guidance below. AI must not invent observations or unsupported parameters; you remain responsible for its output. If you do not use AI, state this and show your estimation method.

\newpage
\section*{Choose 1 Option: model guidance for Task 3}
These are models for \textbf{observations $X$}. No calculus is required. Density has units reciprocal to those of $X$; probabilities are areas, not density heights.

\textbf{A. Linear density on a finite interval (C1--C2, C5).}
Use data with a roughly steadily increasing or decreasing histogram over a plausible finite interval. Choose $a<b$ enclosing all retained observations. Use documented physical bounds if known and suitable; otherwise use $a=\min x_i$, $b=\max x_i$ as simple \textbf{estimated endpoints}. Sample extremes are not known population or physical bounds: new observations could fall outside them. State this limitation. Let $d=b-a$ and choose the direction from the histogram:
\[
f(x)=\begin{cases}2(x-a)/d^2,&a\le x\le b,\\0,&\text{otherwise}\end{cases}
\quad\text{(increasing)},\qquad
f(x)=\begin{cases}2(b-x)/d^2,&a\le x\le b,\\0,&\text{otherwise}\end{cases}
\quad\text{(decreasing)}.
\]
These simple triangular forms have support $[a,b]$ and peak height $2/d$. Verify nonnegativity and triangle area $\tfrac12 d(2/d)=1$. Their model means are $a+2d/3$ (increasing) and $a+d/3$ (decreasing); both have variance $d^2/18$ and SD $d/\sqrt{18}$. Compare with your sample statistics; do not force this model onto an incompatible histogram.

\textbf{B. Exponential waiting times (C1, C4--C5).}
Use nonnegative waiting times with many short waits and fewer long waits, where an approximately constant arrival rate and independent successive waits are plausible. Estimate $\hat\lambda=1/\bar{x}$ for $\bar{x}>0$; the rate has units of inverse time. The fitted density and support are
\[
f(x)=\begin{cases}\hat\lambda e^{-\hat\lambda x},&x\ge0,\\0,&x<0,\end{cases}
\qquad X\in[0,\infty).
\]
The model mean and SD are both $1/\hat\lambda$, and variance is $1/\hat\lambda^2$. Check whether $s$ is reasonably close to $\bar{x}$ and discuss changing rates or dependent waits if relevant.

\textbf{C. Normal measurements (C1, C5; normal methods from C9).}
Use roughly symmetric, unimodal, bell-shaped measurements without severe outliers. Set $\hat\mu=\bar{x}$ and $\hat\sigma=s>0$ as simple estimates, both in the units of $X$. Write $X\sim N(\hat\mu,\hat\sigma^2)$ with
\[
f(x)=\frac{1}{\hat\sigma\sqrt{2\pi}}
\exp\!\left[-\frac{(x-\hat\mu)^2}{2\hat\sigma^2}\right],
\qquad -\infty<x<\infty.
\]
The model mean, variance and SD are $\hat\mu$, $\hat\sigma^2$ and $\hat\sigma$. For physically bounded measurements, discuss whether the normal model gives appreciable probability to impossible values. An unknown population SD does \textbf{not} make raw observations follow Student's $t$. Here $s$ describes observation spread; $s/\sqrt n$ is the standard error of a sample mean. No $t$ inference is needed in this investigation.

\section*{4. Calculate quantiles and model probabilities}
A \textbf{quantile} $q_p$ is a cutoff with $P(X\le q_p)=p$ under your fitted continuous model. Find the \textbf{numerical values and units} of $q_{0.10}$ and $q_{0.90}$. Show parameter substitutions and the steps used to obtain them:
\begin{itemize}
\item \textbf{Linear:} use triangle areas. For $a\le x\le b$, the area to the left is $(x-a)^2/d^2$ (increasing) or $1-(b-x)^2/d^2$ (decreasing). Equating it to $p$ gives $q_p=a+d\sqrt p$ or $q_p=b-d\sqrt{1-p}$, respectively.
\item \textbf{Exponential:} use $P(X\le x)=1-e^{-\hat\lambda x}$ for $x\ge0$, so $q_p=-\ln(1-p)/\hat\lambda$.
\item \textbf{Normal:} standardise with $Z=(X-\hat\mu)/\hat\sigma$. Use normal tables or a calculator to find $P(Z\le z_p)=p$, then $q_p=\hat\mu+\hat\sigma z_p$; $z_{0.10}\approx-1.2816$, $z_{0.90}\approx1.2816$.
\end{itemize}
Calculate the three events
\[
P(X<q_{0.10})=0.10,\qquad
P(q_{0.10}\le X\le q_{0.90})=0.80,\qquad
P(X>q_{0.90})=0.10.
\]
These probabilities hold \textbf{by construction}; you must still show how the numerical cutoffs were obtained. Single endpoints have probability zero in a continuous model.

Choose one additional \textbf{meaningful numerical threshold $c$} from your question, such as an operational limit, and calculate \textbf{$P(X>c)$}. Explain its meaning; do not choose $c$ simply as another model quantile. Use geometric areas, $e^{-\hat\lambda c}$ for exponential $c\ge0$, or normal standardisation and tables/calculator. Respect the support when a threshold lies outside it. Keep sufficient precision in cutoffs and state the numerical boundaries used for comparison.

\newpage
\section*{5. Compare empirical and model probabilities}
For \textbf{all four events in Task 4}, use \textbf{exactly the same numerical boundaries and inequalities} in the dataset and fitted model. Count from observations, not bins crossing the boundaries. If you round cutoffs before counting, evaluate the model at those same cutoffs; its quantile probabilities may then be only approximately $0.10$, $0.80$, $0.10$.

Count the observations satisfying each event and calculate
\[
p_{\mathrm{empirical}}=\frac{\text{event count}}{n},\qquad
E_A=|p_{\mathrm{empirical}}-p_{\mathrm{model}}|,\qquad
E_{\%}=\frac{100E_A}{|p_{\mathrm{model}}|}.
\]
Percentage error is \textbf{undefined if $p_{\mathrm{model}}=0$}. A probability difference of $0.03$ equals \textbf{3 percentage points}, not necessarily $3\%$ relative error. The model probability is the reference for percentage error.

Complete this table with \textbf{numerical boundaries}, decimal probabilities, and percentage errors labelled $\%$.
\begin{center}
\small
\setlength{\tabcolsep}{3pt}
\renewcommand{\arraystretch}{1.2}
\begin{tabularx}{\textwidth}{|>{\raggedright\arraybackslash}X|C{1.5cm}|C{1.5cm}|C{2.2cm}|C{2.2cm}|C{1.9cm}|C{2.2cm}|}
\hline
\textbf{Event} & \textbf{Event count} & \textbf{Sample size $n$} & \textbf{Empirical probability} & \textbf{Model probability} & \textbf{Absolute error} & \textbf{Percentage error}\\\hline
$X<q_{0.10}$ & & & & & & \\\hline
$q_{0.10}\le X\le q_{0.90}$ & & & & & & \\\hline
$X>q_{0.90}$ & & & & & & \\\hline
$X>c$ & & & & & & \\\hline
\end{tabularx}
\end{center}

Empirical frequencies need not equal the quantile probabilities exactly. \textbf{Use fitted-model quantiles, not sample percentiles whose resulting proportions would not provide independent evidence of fit.}

Connect your work to \textbf{C6} (relative frequencies versus theoretical probabilities) and \textbf{C7} (sample versus model statistics). Use \textbf{C1--C5} for supports, density areas, exponential probabilities and moments where relevant, and \textbf{C9} normal standardisation for a normal model. No separate C8 $t$ exercise is required.

\section*{6. Interpret the results}
Write one short conclusion that \textbf{answers your question using at least one calculated probability}. Identify the event with the \textbf{largest absolute error} and explain one plausible reason, such as sampling variation, changing conditions or model shape. Acknowledge that fitting and checking on the \textbf{same dataset is a descriptive assessment}, not proof of accurate predictions for new data.

\section*{Required Report Structure and Submission}
Aim for \textbf{3--4 pages} with a \textbf{title, full name and grade/class}. Use clear headings and connected explanations: \textbf{Introduction and methodology} (question, context, data preparation and links); \textbf{Results and discussion} (statistics, histogram, fitted model/parameters, AI disclosure, quantiles, calculations and comparison table); \textbf{Conclusion and references} (answer, limitations and citations for data and other sources/tools).

Write in \textbf{LaTeX in Overleaf}; submit \textbf{one compiled PDF through Open}. Use \verb|\href{URL}{descriptive text}| with \texttt{hyperref}. All hyperlinks must work in the compiled PDF.

\section*{C10 Assessment: 6 Points}
\textbf{Passing threshold: 5 out of 6 points.}
\begin{center}
\small
\begin{tabularx}{\textwidth}{|>{\raggedright\arraybackslash}X|c|}
\hline
\textbf{Evidence assessed} & \textbf{Points}\\\hline
Question, variable, 30+ real observations, preparation and accessible required data links. & 1\\\hline
Sample statistics with procedures and units; labelled histogram. & 1\\\hline
Suitable density/support, parameter method, checks, limitation and transparent AI disclosure. & 1\\\hline
Two quantiles with working; three quantile probabilities and one contextual probability. & 1\\\hline
Four counts and relative frequencies, matching boundaries, table and both error measures. & 1\\\hline
Evidence-based conclusion, explained mismatch, same-data limitation and clear LaTeX report. & 1\\\hline
\textbf{Total} & \textbf{6}\\\hline
\end{tabularx}
\end{center}

\textbf{Final checklist:} identification and question; accessible source/exact-data links; statistics and histogram; fitted density/support and AI disclosure; quantiles, four events and errors; conclusion and limitations. Open the PDF, test hyperlinks and check readability before uploading to \textbf{Open}.

\end{document}
